% Consolidated 33-thermal-radiation from 620defd2149c8dcfd52d29c147aa4708c27ab8f3; original TeX fragments preserved.

% === BEGIN 33-thermal-radiation-246-257.tex blob c1334462c7481a2d2e0d74ca81ba703f2f71608d ===
% Irodov Part V §5.7, original problems 246--257.
% Source Arihant PDF printed pp.227--228, answers printed pp.337.
% The original printed problem 254 has omega^2, but answer uses omega^3.
% Both literal and physically corrected interpretations are documented.

\thermalopticsproblem{246}{227}
利用维恩的热辐射普遍标度律
$u_\omega(\omega,T)=\omega^3F(\omega/T)$，
证明：(a) 最概然辐射角频率与 $T$ 成正比；
(b) 单位角频率谱能量密度的最大值与 $T^3$ 成正比；
(c) 黑体总辐出度与 $T^4$ 成正比。
\exsolution
令 $x=\omega/T$，
则 $\omega=xT$。
谱能量密度可写成
\[
 u_\omega=T^3x^3F(x).
\]
对固定温度求峰值，
最大值对应的 $x=x_m$ 是函数 $x^3F(x)$
的温度无关极大点，故
\[
 \boxed{\omega_m=x_mT\propto T},\qquad
 \boxed{(u_\omega)_{\max}\propto T^3}.
\]
总能量密度
\[
 u(T)=\int_0^\infty u_\omega\,d\omega
 =T^4\int_0^\infty x^3F(x)\,dx.
\]
对各向同性平衡辐射，
总辐出度 $M_e=cu/4$，因此
\[
 \boxed{M_e\propto T^4}.
\]
\excheck 此处假设积分收敛，
且形状函数 $F$ 本身不显含温度；无需先采用普朗克公式。

\thermalopticsproblem{247}{227}
两个加热黑体中，一个温度为
$T_1=2500\,\mathrm K$。
另一个黑体的辐出度峰值对应波长比第一个大
$\Delta\lambda=0.50\,\mu\mathrm m$。
求第二个黑体的温度。
\exsolution
维恩位移律
$\lambda_m T=b$，取
$b=2.898\times10^{-3}\,\mathrm{m\,K}$。
由
$b/T_2=b/T_1+\Delta\lambda$，
可得
\[
 \boxed{T_2=
 \frac{bT_1}{b+T_1\Delta\lambda}}.
\]
代入
$\Delta\lambda=5.0\times10^{-7}\,\mathrm m$，
\[
 \boxed{T_2\simeq1.75\times10^3\,\mathrm K}.
\]
\excheck 峰值波长变长，温度应降低；
本题 $T_2<T_1$ 满足这一极限检查。

\thermalopticsproblem{248}{227}
某黑体的总辐出度为
$M_e=3.0\,\mathrm{W/cm^2}$。
求该黑体光谱辐出度最大值对应的波长。
\exsolution
首先换算为
$M_e=3.0\times10^4\,\mathrm{W/m^2}$。
斯特藩--玻尔兹曼定律给
$M_e=\sigma T^4$，所以
$T=(M_e/\sigma)^{1/4}\simeq853\,\mathrm K$。
再由维恩位移律
\[
 \boxed{\lambda_m=
 b\left(\frac{\sigma}{M_e}\right)^{1/4}
 \simeq3.40\,\mu\mathrm m}.
\]
\excheck 用的是对\emph{波长谱}的峰值所对应的
位移常数 $b$，不能改用频率谱峰值的不同常数。

\thermalopticsproblem{249}{227}
太阳的辐射谱可近似看作黑体谱，
其波长谱峰值位于
$\lambda_m=0.48\,\mu\mathrm m$。
估算太阳每秒因辐射失去的质量，
以及按此恒定辐射率失去总质量的 $1\%$
所需时间。
\exsolution
太阳表面有效温度由维恩位移律得到
\[
 T\simeq\frac{b}{0.48\times10^{-6}}
 \simeq6.04\times10^3\,\mathrm K.
\]
取太阳半径
$R_\odot\simeq6.96\times10^8\,\mathrm m$，
光度（总辐射功率）
$L=4\pi R_\odot^2\sigma T^4
\simeq4.6\times10^{26}\,\mathrm W$。
由质能等价
$dM/dt=-L/c^2$，得
\[
 \boxed{\left|\frac{dM}{dt}\right|
 \simeq5.1\times10^9\,\mathrm{kg/s}}.
\]
若以
$M_\odot\simeq1.99\times10^{30}\,\mathrm{kg}$
且忽略太阳光度、半径演化，
对应 $1\%$ 质量损失时间为
\[
 \boxed{\Delta t=
 \frac{0.01M_\odot c^2}{L}
 \simeq1.2\times10^{11}\,\mathrm{yr}}.
\]
\excheck 这是\emph{固定光度的形式外推}；
真实太阳演化不允许把该数值当成寿命预测。

\thermalopticsproblem{250}{227}
完全电离的氢等离子体，密度为
$\rho=0.10\,\mathrm{g/cm^3}$。
求热辐射压与等离子体气体动理论压强相等时
的温度 $T$。
取平衡辐射压 $p_{\rm rad}=u/3$，
其中辐射能量密度 $u=4\sigma T^4/c$；
等离子体可视为理想气体。
\exsolution
完全电离氢每个氢核伴随一个自由电子，
忽略电子质量，每单位体积粒子总数
$n_{\rm tot}=2\rho/m_p$
（或者写 $2\rho N_A/M_H$），于是
\[
 p_{\rm gas}=n_{\rm tot}k_BT
 =\frac{2\rho R_g}{M_H}T.
\]
辐射压
$p_{\rm rad}=4\sigma T^4/(3c)$。
令二者相等，约去温度 $T>0$：
\[
 \frac{4\sigma T^3}{3c}
 =\frac{2\rho R_g}{M_H},
 \qquad
 \boxed{T=\left(
 \frac{3c\rho R_g}{2\sigma M_H}
 \right)^{1/3}}.
\]
代
$\rho=100\,\mathrm{kg/m^3}$、
$M_H\simeq1.008\times10^{-3}\,\mathrm{kg/mol}$，
\[
 \boxed{T\simeq1.9\times10^7\,\mathrm K}.
\]
\excheck 氢的\emph{质子和电子各贡献一份气体压强}，
因子2不可漏掉；核反应与电子简并均未计入。

\thermalopticsproblem{251}{227}
一颗直径 $d=1.2\,\mathrm{cm}$ 的铜球放在
内壁绝对零度且抽成真空的容器中。
初温 $T_0=300\,\mathrm K$，视铜球表面为理想黑体，
求降到初温的 $1/\eta$（$\eta=2.0$）所需时间。
\exsolution
设铜的质量密度为 $\rho$、比热容为 $c_p$，
铜球质量
$m=\rho \pi d^3/6$，
辐射面积 $A=\pi d^2$。
忽略物性温度变化，能量平衡
\[
 mc_p\frac{dT}{dt}=-\sigma AT^4.
\]
分离变量从 $T_0$ 积分到
$T_f=T_0/\eta$：
\[
 \int_{T_f}^{T_0}\frac{dT}{T^4}
 =-\frac{\sigma A}{mc_p}\int_{t}^{0}\!dt,
\]
等价地
$\tfrac13(T_f^{-3}-T_0^{-3})
=\sigma At/(mc_p)$。
所以
\[
 \boxed{t=\frac{c_p\rho d}{18\sigma T_0^3}
 (\eta^3-1)}.
\]
取 $\rho\simeq8.96\times10^3\,\mathrm{kg/m^3}$、
$c_p\simeq385\,\mathrm{J/(kg\,K)}$，
\[
 \boxed{t\simeq1.05\times10^4\,\mathrm s
 \simeq2.9\,\mathrm h\approx3\,\mathrm h}.
\]
\excheck 原题理想黑体与真空壁为0K的重要条件
使净辐射功率恰为 $\sigma AT^4$；
若环境温度不为0，应减去环境辐射项。

\thermalopticsproblem{252}{227--228}
两个黑体腔各有一个直径
$d=1.0\,\mathrm{cm}$ 的小圆孔，
两孔彼此同轴、相向放置，相距
$l=10\,\mathrm{cm}$（原图5.39），
腔外表面理想反射。
第一个腔维持
$T_1=1700\,\mathrm K$。
求第二个腔达到稳态时的温度。
考虑黑体孔的余弦发射规律。
\exsolution
黑体孔的总辐出度为 $\sigma T^4$，
近似平行的小圆孔之间的辐射视角因子为
\[
 F_{1\to2}\simeq\frac{A_2}{\pi l^2}
 =\frac{\pi d^2/4}{\pi l^2}
 =\frac{d^2}{4l^2}.
\]
故第二腔每秒通过孔2吸收的来自孔1功率为
$\sigma T_1^4 A_1 F_{1\to2}$。
稳态时孔2向外发出
$\sigma T_2^4 A_2$，
由于两孔面积相同 $A_1=A_2$，所以
\[
 T_2^4=T_1^4\frac{d^2}{4l^2}.
\]
开四次方得
\[
 \boxed{T_2=T_1\sqrt{\frac d{2l}}
 =1700\sqrt{\frac{1.0}{20}}
 \simeq380\,\mathrm K\approx0.4\,\mathrm{kK}}.
\]
\excheck 应开四次方，故温度是
$T_1\sqrt{d/(2l)}$；
若在书末纯文本中见到 $T_1d/(2l)$，
属于根号在文本提取时丢失。

\thermalopticsproblem{253}{228}
一个容积 $V=1.0\,\mathrm L$ 的黑体腔，
充满温度 $T=1000\,\mathrm K$ 的平衡热辐射。
求 (a) 定容热容 $C_V$；
(b) 辐射的熵 $S$。
\exsolution
平衡辐射能量密度
$u=4\sigma T^4/c$，
故
$U=uV=4\sigma VT^4/c$，
定容热容
\[
 \boxed{C_V=\left(\frac{\partial U}{\partial T}\right)_V
 =\frac{16\sigma VT^3}{c}}.
\]
黑体光子气满足
$p=u/3$，并有
$TS=U+pV=4U/3$（化学势为0），所以
\[
 \boxed{S=\frac{16\sigma VT^3}{3c}}.
\]
以 $V=10^{-3}\,\mathrm{m^3}$ 代入，
\[
 \boxed{C_V\simeq3.0\,\mathrm{nJ/K}},\qquad
 \boxed{S\simeq1.0\,\mathrm{nJ/K}}.
\]
\excheck 熵与热容在本问题分别随
$T^3$ 增长，且 $S=C_V/3$。

\thermalopticsproblem{254}{228}
已知热辐射维恩角频率谱
\[
 \boxed{u_\omega(\omega,T)=A\omega^3e^{-a\omega/T}},
 \quad a=7.64\,\mathrm{ps\,K},\quad T=2000\,\mathrm K.
\]
分别求(a) 单位角频率谱密度的峰值位置
$\omega_{\max,\omega}$；
(b) 单位波长谱密度的峰值位置
$\lambda_{\max,\lambda}$。
\textbf{版本说明：}这一 $\omega^3$ 形式与标准
Planck 公式的高频极限及可查 Mir 1988 英文题面一致；
但本项目此前记录的 Arihant 扫描底本
印成 $\omega^2$，故另列其字面解，保留来源差异。
\exsolution
设 $D=a/T$。
对标准物理版本，
\[
 \frac{d}{d\omega}
 \bigl(\omega^3e^{-D\omega}\bigr)
 =\omega^2e^{-D\omega}(3-D\omega).
\]
频率谱峰值因此为
\[
 \boxed{\omega_{\max,\omega}=\frac{3T}{a}
  \simeq7.85\times10^{14}\,\mathrm{rad/s}}.
\]
波长谱极大不能直接将该峰值用
$\lambda=2\pi c/\omega_{\max,\omega}$ 代换：
谱密度变换必须使用变量雅可比。
由 $\omega=2\pi c/\lambda$，
$\left|d\omega/d\lambda\right|=2\pi c/\lambda^2$，
\[
 u_\lambda(\lambda,T)=
 u_\omega(2\pi c/\lambda,T)
 \left|\frac{d\omega}{d\lambda}\right|
 \propto\lambda^{-5}
 \exp\!\left(-\frac{2\pi ca}{T\lambda}\right).
\]
极值条件
\[
 \frac{d}{d\lambda}\ln u_\lambda
 =-\frac5\lambda+\frac{2\pi ca}{T\lambda^2}=0
\]
给出
\[
 \boxed{\lambda_{\max,\lambda}
  =\frac{2\pi ca}{5T}
  \simeq1.44\,\mathrm{\mu m}}.
\]
\textbf{按此前记录扫描底本 $\omega^2$
的独立字面分支：}
若题设实际是
$u_\omega=A\omega^2e^{-a\omega/T}$，
则角频率峰在
$\boxed{\omega_{\max,\omega}^{(2)}=2T/a
 \simeq5.24\times10^{14}\,\mathrm{rad/s}}$；
且由于对应
$u_\lambda\propto\lambda^{-4}
 \exp[-2\pi ca/(T\lambda)]$，
\[
 \boxed{\lambda_{\max,\lambda}^{(2)}
 =\frac{2\pi ca}{4T}
 =\frac{\pi ca}{2T}
 \simeq1.80\,\mathrm{\mu m}}.
\]
\excheck
两谱的峰位不是同一个光子、也非互为
$\lambda=2\pi c/\omega$ 的简单倒数；
标准形式的峰位满足
$\lambda_{\max,\lambda}\omega_{\max,\omega}
 =\frac35(2\pi c)$，
打印为 $\omega^2$ 时则为
$\frac12(2\pi c)$。
标准维恩系数 $a=\hbar/k_B$，
因此从 Planck 高频近似必须得到 $\omega^3$。
不同版本对照见
\texttt{sinhafoundation.org/physics/problems/5-254/}，
不能把两个谱模型的频率/波长极值混用。

\thermalopticsproblem{255}{228}
从普朗克黑体辐射公式推导角频率谱能量密度
$u_\omega$ 的两个渐近形式：
(a) $\hbar\omega\ll k_BT$（瑞利--金斯极限）；
(b) $\hbar\omega\gg k_BT$（维恩极限）。
\exsolution
普朗克公式
\[
 u_\omega=
 \frac{\hbar\omega^3}{\pi^2c^3}
 \frac{1}{e^{\hbar\omega/(k_BT)}-1}.
\]
记 $x=\hbar\omega/(k_BT)>0$。
(a) $x\ll1$ 时
$e^x-1=x+O(x^2)$，
因而
\[
 \boxed{u_\omega\simeq
 \frac{k_BT\,\omega^2}{\pi^2c^3}}.
\]
(b) $x\gg1$ 时
$e^x-1\simeq e^x$，
故
\[
 \boxed{u_\omega\simeq
 \frac{\hbar\omega^3}{\pi^2c^3}
 e^{-\hbar\omega/(k_BT)}}.
\]
\excheck 两式都具有单位
$\mathrm{J\,m^{-3}\,(rad/s)^{-1}}$，
且后者的指数必须无量纲。

\thermalopticsproblem{256}{228}
将普朗克公式的角频率谱能量密度
$u_\omega$ 改写为普通频率 $\nu$
以及真空波长 $\lambda$ 的谱能量密度
$u_\nu$、$u_\lambda$。
\exsolution
由于谱元的能量必须相等，
\[
 u_\omega\,d\omega=u_\nu\,d\nu
 =-u_\lambda\,d\lambda
\quad\text{（最后一个等式的负号来自变量反向）}.
\]
普通频率与角频率的关系
$\omega=2\pi\nu$、
$d\omega/d\nu=2\pi$。
令 $h=2\pi\hbar$，得到
\[
 \boxed{u_\nu(\nu,T)=
 \frac{8\pi h\nu^3}{c^3}
 \frac{1}{e^{h\nu/(k_BT)}-1}}.
\]
再以 $\nu=c/\lambda$，
取绝对雅可比因子
$|d\nu/d\lambda|=c/\lambda^2$，
所以
\[
 \boxed{u_\lambda(\lambda,T)=
 \frac{8\pi hc}{\lambda^5}
 \frac{1}{e^{hc/(\lambda k_BT)}-1}}.
\]
\excheck $u_\nu$ 与 $u_\lambda$ 的数值单位不同；
例如 $u_\lambda$ 是单位\emph{波长}下的能量密度，
不能不带雅可比因子直接用 $\nu=c/\lambda$ 代入。

\thermalopticsproblem{257}{228}
黑体温度为 $T=3000\,\mathrm K$。
求其\emph{单位面积}在接近
波长谱能量密度最大值处的
$\Delta\lambda=1.0\,\mathrm{nm}$ 窄波段内
所辐射的功率。
\exsolution
波长谱辐出度
$M_\lambda=(c/4)u_\lambda$，
由普朗克公式得
\[
 M_\lambda=
 \frac{2\pi hc^2}{\lambda^5}
 \frac{1}{e^{hc/(\lambda k_BT)}-1}.
\]
波长谱峰值位于
$\lambda_m=b/T$，
取 $b=2.89777\times10^{-3}\,\mathrm{m\,K}$，
有
$\lambda_m\simeq9.659\times10^{-7}\,\mathrm m$。
因为 $\Delta\lambda\ll\lambda_m$，
窄波段辐出度近似为
\[
 \Delta M\simeq M_\lambda(\lambda_m)\Delta\lambda.
\]
将 $T=3000\,\mathrm K$、
$\Delta\lambda=10^{-9}\,\mathrm m$
及精确常数 $h,c,k_B$ 代入，
得
\[
 \boxed{\Delta M\simeq
 3.13\times10^3\,\mathrm{W/m^2}
 =0.313\,\mathrm{W/cm^2}}.
\]
\excheck 相比黑体总辐出度
$\sigma T^4\simeq4.59\times10^6\,\mathrm{W/m^2}$，
这个1nm带宽只占很小部分，数量级合理。
% === END 33-thermal-radiation-246-257.tex ===

% === BEGIN 33-thermal-radiation-258-268.tex blob 01e1f46c580ce87770d01d971aafa79c69424700 ===
% Irodov Part V §5.7, originals 258--268 (printed pp.228--229).
% Source: user-supplied Arihant PDF, actual pp.237--238.
% Newly worked derivations; figures 5.40 and relevant boundary conventions
% remain attached to source, no third-party academic review is asserted.
\thermalopticsproblem{258}{228--229}
原图5.40给出黑体辐射累积功率份额 $y(x)$，其中
$x=\lambda/\lambda_m$，$\lambda_m$ 为波长谱辐出度峰值波长。
据此求：(a) $T=3700\,\mathrm K$ 时将总辐射功率平分的波长；
(b) $T=5000\,\mathrm K$ 时落在可见光
$0.40$--$0.76\,\mu\mathrm m$ 内的功率份额；
(c) 温度从 $3000\,\mathrm K$ 升为 $5000\,\mathrm K$ 时
$\lambda>0.76\,\mu\mathrm m$ 部分的辐射功率增加倍数。
\exsolution
对黑体波长谱使用普朗克公式
\[
 M_\lambda(T)=\frac{2\pi hc^2}{\lambda^5}
 \frac1{\exp\!\left(hc/(\lambda k_{\mathrm B}T)\right)-1}.
\]
令 $b=2.897771955\times10^{-3}\,\mathrm{m\,K}$，
$\lambda_m=b/T$，累积份额写作
\[
 Y(x)=\frac{\displaystyle\int_0^{x\lambda_m}M_\lambda(T)\,d\lambda}
 {\sigma T^4}.
\]
换元 $\lambda=x\lambda_m$ 可知 $Y$ 与 $T$ 无关，
仅依赖 $x$；可由原图近似读数，或作上述定积分。
(a) 数值求解 $Y(x_{1/2})=1/2$ 得
$x_{1/2}\simeq1.41738$，于是
\[
 \boxed{\lambda_{1/2}=\frac{bx_{1/2}}{3700}
 \simeq1.11\,\mu\mathrm m}.
\]
(b) 令 $x_1=0.40T/b\simeq0.6902$、
$x_2=0.76T/b\simeq1.3114$（此处微米与 $b$ 一致换算），
\[
 \boxed{f_{\mathrm{vis}}=Y(x_2)-Y(x_1)\simeq0.377}.
\]
(c) 记 $x_T=0.76\,\mu\mathrm m\,T/b$。
波长超过截止值的辐出度为
$\sigma T^4[1-Y(x_T)]$，所以
\[
 \boxed{\eta=\left(\frac{5000}{3000}\right)^4
 \frac{1-Y(x_{5000})}{1-Y(x_{3000})}
 \simeq4.86}.
\]
\excheck (a) 是\emph{功率平分波长}，不是谱峰波长；
(b) 为功率比例而非可见光的波长区间比例；
(c) 不可只用 $T^4$，还须修正所选波段份额。

\thermalopticsproblem{259}{229}
从普朗克辐射谱推导温度为 $T$ 的黑体腔内，
在角频率间隔 $(\omega,\omega+d\omega)$
及波长间隔 $(\lambda,\lambda+d\lambda)$ 中，
每立方厘米的平均光子数。
\exsolution
每单位体积、单位角频率内的辐射能量密度为
\[
 u_\omega=\frac{\hbar\omega^3}{\pi^2c^3}
 \frac1{e^{\hbar\omega/(k_{\mathrm B}T)}-1}.
\]
每光子能量为 $\hbar\omega$，故光子数谱密度
$n_\omega=u_\omega/(\hbar\omega)$；
利用 $\omega=2\pi c/\lambda$ 和
$|d\omega/d\lambda|=2\pi c/\lambda^2$，得到
\[
 \boxed{dN_\omega=
 10^{-6}\frac{\omega^2\,d\omega}{\pi^2c^3
 [e^{\hbar\omega/(k_{\mathrm B}T)}-1]}}
 \quad\text{每}\;\mathrm{cm^3},
\]
\[
 \boxed{dN_\lambda=
 10^{-6}\frac{8\pi\,d\lambda}{\lambda^4
 [e^{hc/(\lambda k_{\mathrm B}T)}-1]}}
 \quad\text{每}\;\mathrm{cm^3},
\]
其中 $\omega,\lambda$ 及区间宽度使用 SI 单位。
\excheck 角频率谱转波长谱必须引入
$|d\omega/d\lambda|$，而且 $1\,\mathrm{cm^3}=10^{-6}\,\mathrm{m^3}$。

\thermalopticsproblem{260}{229}
功率 $P=10\,\mathrm W$ 的各向同性点光源
发射波长 $\lambda=589\,\mathrm{nm}$ 的单色光。
求：(a) 距光源 $r=2.0\,\mathrm m$ 处
平均光子流密度；(b) 光子平均数密度降为
$n=100\,\mathrm{cm^{-3}}$ 的距离。
\exsolution
单个光子能量 $\varepsilon=hc/\lambda$，
各向同性源在半径 $r$ 球面上能流密度
$P/(4\pi r^2)$，所以
\[
 \boxed{j_N(r)=\frac{P\lambda}{4\pi r^2hc}}.
\]
代入 $r=2.0\,\mathrm m$，
\[
 \boxed{j_N\simeq5.90\times10^{17}\,
 \mathrm{m^{-2}s^{-1}}}.
\]
自由空间内沿径向传播的光子数密度为 $n=j_N/c$。
将 $100\,\mathrm{cm^{-3}}=10^8\,\mathrm{m^{-3}}$ 代入，
\[
 \boxed{r=\sqrt{\frac{P\lambda}{4\pi hc^2n}}
 \simeq8.87\,\mathrm m}.
\]
\excheck 题目(b)给的是单位\emph{体积}的光子数，
不能误用(a)的单位面积光子通量。

\thermalopticsproblem{261}{229}
从光子观点证明：平行光束单位时间传给
完全吸收靶面的动量，仅由光束总能流率
$\Phi_e$ 决定，与辐射光谱组成无关。
\exsolution
任意频率为 $\nu_i$ 的光子能量为
$\varepsilon_i=h\nu_i$，动量模为
$p_i=\varepsilon_i/c$。
若单位时间入射的第 $i$ 类光子数为 $\dot N_i$，
且所有光子均被吸收，单位时间传递的动量
\[
 \frac{dp}{dt}=\sum_i\dot N_i p_i
 =\frac1c\sum_i\dot N_i\varepsilon_i.
\]
括号内的能量总和恰为入射功率 $\Phi_e$，因此
\[
 \boxed{\frac{dp}{dt}=\frac{\Phi_e}{c}}.
\]
\excheck 若靶面理想镜面反射且正入射，
动量传递率改成 $2\Phi_e/c$；
本题所给书末关系对应全吸收条件。

\thermalopticsproblem{262}{229}
激光脉冲历时 $\tau=0.13\,\mathrm{ms}$，
能量 $E=10\,\mathrm J$，
聚焦到直径 $d=10\,\mu\mathrm m$ 的圆形光斑，
垂直入射表面，表面反射率 $\rho=0.50$。
求平均光压。
\exsolution
对于不透光表面，单位入射能量的法向动量转移
为 $(1+\rho)/c$（吸收部分贡献1倍，
反射部分贡献2倍）。
脉冲平均功率 $E/\tau$，光斑面积
$A=\pi d^2/4$，故
\[
 \boxed{p=\frac{(1+\rho)E}{Ac\tau}
 =\frac{4(1+\rho)E}{\pi d^2c\tau}}.
\]
数值为
\[
 \boxed{p\simeq4.90\times10^6\,\mathrm{Pa}
 \simeq48\,\mathrm{atm}\approx50\,\mathrm{atm}}.
\]
\excheck 必须将 $0.13\,\mathrm{ms}$ 换成
$1.3\times10^{-4}\,\mathrm s$；该值是脉冲持续期间的平均压强。

\thermalopticsproblem{263}{229}
一束能量 $E=7.5\,\mathrm J$ 的短光脉冲
以入射角 $\theta=30^\circ$ 照射不透光反射板，
反射率 $\rho=0.60$。用光子观点求
传给板的总动量\emph{模}。
\exsolution
取 $\hat{\boldsymbol n}$ 为板的法线，
$\hat{\boldsymbol t}$ 为切线。
入射总光子动量大小 $E/c$。
镜面反射部分的法向动量反号、切向不变；
吸收部分的动量全部交给板。
于是板所获动量分量为
\[
 \Delta p_n=\frac Ec(1+\rho)\cos\theta,\qquad
 \Delta p_t=\frac Ec(1-\rho)\sin\theta.
\]
因此
\[
 \boxed{\Delta p=\frac Ec
 \sqrt{1+\rho^2+2\rho\cos(2\theta)}
 \simeq3.50\times10^{-8}\,\mathrm{kg\,m/s}}.
\]
\excheck 光子束有非零切向动量损失；
只取法向分量无法回答“总动量模”。
该数值也可写为 $35\,\mathrm{nN\,s}$。

\thermalopticsproblem{264}{229}
强度 $I=0.20\,\mathrm{W/cm^2}$ 的平面光波
以入射角 $\theta=45^\circ$ 照射反射率
$\rho=0.80$ 的不透光平面镜。
用光子观点求单位\emph{镜面面积}所受法向光压。
\exsolution
$I$ 是垂直于光束方向单位面积的能流密度；
单位镜面面积每秒入射能量为 $I\cos\theta$。
每单位入射能量的法向动量传递为
$(1+\rho)\cos\theta/c$，所以
\[
 \boxed{p=\frac Ic(1+\rho)\cos^2\theta}.
\]
取 $I=2000\,\mathrm{W/m^2}$ 得
\[
 \boxed{p\simeq6.0\times10^{-6}\,\mathrm{Pa}
 =0.60\,\mathrm{nN/cm^2}}.
\]
\excheck 两个 $\cos\theta$ 分别来自
\emph{斜入射通量投影}与\emph{法向动量投影}。

\thermalopticsproblem{265}{229}
强度 $I=0.70\,\mathrm{W/cm^2}$ 的平行光
照射半径 $R=5.0\,\mathrm{cm}$ 的理想镜面球，
从光子观点求球所受的合力。
\exsolution
令 $\theta$ 为球面法线与入射光方向的锐角。
镜面反射造成局部法向压强
$p(\theta)=2I\cos^2\theta/c$，
光沿传播方向的力分量再乘 $\cos\theta$。
照明半球面积元为
$dA=2\pi R^2\sin\theta\,d\theta$，横向力对称抵消。
沿入射传播方向合力
\[
 F=\frac{4\pi IR^2}{c}
 \int_0^{\pi/2}\cos^3\theta\sin\theta\,d\theta
 =\boxed{\frac{\pi R^2I}{c}}.
\]
代入参数，
\[
 \boxed{F\simeq1.83\times10^{-7}\,\mathrm N
 =0.18\,\mu\mathrm N}.
\]
\excheck 球面的积分与平面镜不同；
不能简单把球投影面积乘 $2I/c$。

\thermalopticsproblem{266}{229}
功率为 $P$ 的各向同性点光源位于
理想镜面圆盘的轴上，
离盘心距离 $h=\eta R$，
其中 $R$ 为盘半径。求光源对圆盘的轴向作用力。
\exsolution
令盘上距中心径向距离为 $r$ 的面积元
$dA=2\pi r\,dr$，光源到该元距离
$s=\sqrt{h^2+r^2}$，入射角满足 $\cos\theta=h/s$。
各向同性光强为 $P/(4\pi s^2)$；
镜面反射的法向压力为
$2[P/(4\pi s^2)]\cos^2\theta/c$，所以
\[
 F=\frac{P}{c}\int_0^R
 \frac{h^2r}{(h^2+r^2)^2}\,dr
 =\frac{P}{2c}\frac{R^2}{h^2+R^2}.
\]
由 $h=\eta R$，
\[
 \boxed{F=\frac{P}{2c(1+\eta^2)}}.
\]
\excheck 当 $R\ll h$ 时 $F\propto R^2/h^2$；
面板半径增大时 $F$ 单调增加。

\thermalopticsproblem{267}{229}
实验室系 $K$ 中角频率 $\omega$ 的单个光子
垂直射向以相对论速度 $V$ 迎面运动的理想反射镜。
在忽略镜面反冲速度变化的情况下求反射期间
镜面所获动量：(a) 镜面瞬时静止系；
(b) 实验室系 $K$。
\exsolution
令 $\beta=V/c>0$，光子能量 $E=\hbar\omega$。
对镜面静止系作多普勒变换，
$E'=E\sqrt{(1+\beta)/(1-\beta)}$；
在此系中光子反向而能量不变，所以
\[
 \boxed{\Delta p'=\frac{2E'}c
 =\frac{2\hbar\omega}{c}
 \sqrt{\frac{1-\beta^2}{(1-\beta)^2}}}.
\]
实验室系中的反射光子能量为
$E_r=E(1+\beta)/(1-\beta)$，
入射与出射动量方向相反，故
\[
 \boxed{\Delta p_K=\frac{E+E_r}{c}
 =\frac{2\hbar\omega}{c(1-\beta)}}.
\]
\excheck $\beta\to0$ 时两式均趋于
$2\hbar\omega/c$；镜面接近光子时反射光蓝移，
实验室动量转移相应增加。

\thermalopticsproblem{268}{229}
质量 $m=10\,\mathrm{mg}$ 的小理想平面镜，
用长 $\ell=10\,\mathrm{cm}$ 的无质量细线悬挂。
水平垂直向镜面发射能量为 $E=13\,\mathrm J$
的短激光脉冲，求镜面偏转最大角
$\theta$，并说明镜子的动能来自何处。
\exsolution
入射与反射光子的总法向动量差（小反冲近似）
为 $p=2E/c$。
镜子的初动能
$K=p^2/(2m)=2E^2/(mc^2)$，
在摆动最高点转化为重力势能
$mg\ell(1-\cos\theta)$。
所以
\[
 2mg\ell\sin^2\frac\theta2
 =\frac{2E^2}{mc^2},\qquad
 \boxed{\sin\frac\theta2=\frac{E}{mc\sqrt{g\ell}}}.
\]
代入得
\[
 \boxed{\theta\simeq0.50^\circ}.
\]
\excheck 反射光因镜面反冲而发生极微小的红移，
所失光能变成镜子的动能；
不能以“理想反射”误认为无能量转移。
% === END 33-thermal-radiation-258-268.tex ===

% === BEGIN 33-thermal-radiation-269-280.tex blob 3e83acdadf0e666135a3aceb9790752ce54454ab ===
% Irodov Part V §5.7, originals 269--280, printed p.230.
% Source: user-supplied Arihant PDF p.239; work functions from Appendix 22
% (PDF p.378), zinc=3.74 eV, copper=4.47 eV, cesium=1.89 eV, lithium=2.39 eV.

\thermalopticsproblem{269}{230}
一颗质量为 $M$、半径为 $R$ 的恒星表面发出
角频率为 $\omega_0$ 的光子，求远处观测者接收光子的
相对引力频移 $\Delta\omega/\omega_0$。
\exsolution
在恒星外部近似采用史瓦西静态时空。
静态观察者测得的局部频率满足
$\omega(r)\sqrt{1-2GM/(rc^2)}=\mathrm{const}$。
远处的频率因而是
\[
 \omega_\infty=\omega_0\sqrt{1-\frac{2GM}{Rc^2}}.
\]
故
\[
 \boxed{\frac{\Delta\omega}{\omega_0}
 =\sqrt{1-\frac{2GM}{Rc^2}}-1
 \simeq-\frac{GM}{Rc^2}},
\]
后一式用于 $2GM/(Rc^2)\ll1$。
\excheck 光子向外克服引力势，远处频率\emph{降低}。
题目只给 $M,R$，不可臆造具体数值。

\thermalopticsproblem{270}{230}
X 射线管电压升高到原来的
$\eta=1.5$ 倍，连续谱短波限移动
$\Delta\lambda=26\,\mathrm{pm}$。
求原电压 $U_1$。
\exsolution
连续谱短波限由单电子能量全部转为光子给出
$\lambda_{\min}=hc/(eU)$。
原、后波长分别为
$\lambda_1=hc/(eU_1)$、
$\lambda_2=hc/(e\eta U_1)$，
故
\[
 \Delta\lambda=\lambda_1-\lambda_2
 =\frac{hc}{eU_1}\left(1-\frac1\eta\right).
\]
整理得到
\[
 \boxed{U_1=\frac{hc}{e\Delta\lambda}
 \left(1-\frac1\eta\right)
 \simeq15.9\,\mathrm{kV}}.
\]
\excheck 升高电压时短波限向\emph{更短}波长移动；
$\Delta\lambda$ 是减少量的正值。

\thermalopticsproblem{271}{230}
窄 X 射线束照射 NaCl 单晶，
可观察到镜面布拉格反射的最小掠射角为
$\alpha=4.1^\circ$，
晶面距 $d=0.28\,\mathrm{nm}$。
求射线管的加速电压。
\exsolution
满足一级布拉格反射的最短可用波长为
$\lambda_{\min}=2d\sin\alpha$；
它也必须等于连续谱截止波长
$hc/(eU)$。
因此
\[
 \boxed{U=\frac{hc}{2ed\sin\alpha}
 \simeq31.0\,\mathrm{kV}}.
\]
\excheck 所用 $\alpha$ 是相对晶面的\emph{掠射角}，
不是相对法线的入射角；截止处取布拉格一级 $m=1$。

\thermalopticsproblem{272}{230}
X 射线管内电子到达靶前速度为 $v=0.85c$。
求轫致辐射连续谱的最短波长。
\exsolution
电子为相对论粒子，其动能为
$K=(\gamma-1)m_ec^2$，
$\gamma=(1-v^2/c^2)^{-1/2}$。
最短波长来自全部动能转为单光子：
\[
 \boxed{\lambda_{\min}=
 \frac{hc}{(\gamma-1)m_ec^2}}.
\]
用 $\gamma\simeq1.8983$、
$m_ec^2\simeq511\,\mathrm{keV}$，
\[
 \boxed{\lambda_{\min}\simeq2.70\,\mathrm{pm}
 \approx2.8\,\mathrm{pm}}.
\]
\excheck $v=0.85c$ 时不能用非相对论
$m_ev^2/2$ 代替电子动能。

\thermalopticsproblem{273}{230}
求锌的光电效应截止波长，
以及波长 $250\,\mathrm{nm}$ 光照下
锌表面光电子的最大速度。
\exsolution
原书附录22给锌的逸出功
$A_{\mathrm{Zn}}=3.74\,\mathrm{eV}$。
截止处 $hc/\lambda_{\mathrm{th}}=A$，所以
\[
 \boxed{\lambda_{\mathrm{th}}=
 \frac{hc}{A_{\mathrm{Zn}}}\simeq332\,\mathrm{nm}}.
\]
入射光子能量为
$hc/(250\,\mathrm{nm})\simeq4.96\,\mathrm{eV}$，
故最大电子动能
$K_{\max}\simeq1.22\,\mathrm{eV}$，满足
$K_{\max}=m_ev_{\max}^2/2$，所以
\[
 \boxed{v_{\max}=
 \sqrt{\frac{2}{m_e}
 \left(\frac{hc}{250\,\mathrm{nm}}-A_{\mathrm{Zn}}\right)}
 \simeq6.55\times10^5\,\mathrm{m/s}}.
\]
\excheck 两个光子能量均按 eV 一致换算；
锌的截止波长大于入射波长，确实发生光电发射。

\thermalopticsproblem{274}{230}
同一金属分别以
$\lambda_1=0.35\,\mu\mathrm m$、
$\lambda_2=0.54\,\mu\mathrm m$ 的光照射，
两次产生的光电子最大速度之比为
$\eta=v_1/v_2=2.0$。
求该金属逸出功 $A$。
\exsolution
爱因斯坦光电方程给
$\tfrac12m_ev_i^2=hc/\lambda_i-A$。
由动能比为速度比平方，
\[
 \frac{hc/\lambda_1-A}
 {hc/\lambda_2-A}=\eta^2.
\]
解得
\[
 \boxed{A=
 \frac{\eta^2hc/\lambda_2-hc/\lambda_1}
 {\eta^2-1}
 \simeq1.88\,\mathrm{eV}\approx1.9\,\mathrm{eV}}.
\]
\excheck 不能把速度比 $\eta$ 当作能量比，
应使用 $\eta^2=4$。

\thermalopticsproblem{275}{230}
一个远离其他导体的铜球受到
波长 $\lambda=140\,\mathrm{nm}$ 的辐射。
在无漏电、弱场近似下，球最多能被充到多少正电势？
\exsolution
铜的逸出功按原书附录22为
$A_{\mathrm{Cu}}=4.47\,\mathrm{eV}$。
光电子最大动能
$K_{\max}=hc/\lambda-A_{\mathrm{Cu}}$。
随着球失去电子，其正电势增高，
出射电子还须克服 $e\varphi$；
停止发射条件为 $e\varphi_{\max}=K_{\max}$。
故
\[
 \boxed{\varphi_{\max}=
 \frac1e\left(\frac{hc}{\lambda}-A_{\mathrm{Cu}}\right)
 \simeq4.39\,\mathrm V\approx4.4\,\mathrm V}.
\]
\excheck 球在光电发射后带\emph{正电}；
若 $hc/\lambda<A_{\mathrm{Cu}}$，则不能仅靠此光致正充电。

\thermalopticsproblem{276}{230}
一束照射锂表面的调幅电磁辐射，
电场为
$E(t)=a(1+\cos\omega t)\cos\omega_0t$，
其中 $\omega=6.0\times10^{14}\,\mathrm{s^{-1}}$，
$\omega_0=3.60\times10^{15}\,\mathrm{s^{-1}}$。
求逸出光电子的最大动能。
\exsolution
先将调制信号展开：
\[
 E(t)=a\cos\omega_0t+
 \frac a2\cos(\omega_0+\omega)t+
 \frac a2\cos(\omega_0-\omega)t.
\]
三个谱线的角频率为
$\omega_0,\omega_0\pm\omega$，
最大光子能量对应高频边带 $\omega_0+\omega$。
锂逸出功为原书附录22的
$A_{\mathrm{Li}}=2.39\,\mathrm{eV}$，因此
\[
 \boxed{K_{\max}=
 \hbar(\omega_0+\omega)-A_{\mathrm{Li}}
 \simeq0.375\,\mathrm{eV}\approx0.38\,\mathrm{eV}}.
\]
\excheck 最大动能由\emph{最高频的频谱分量}决定，
不是由电场振幅大小决定。

\thermalopticsproblem{277}{230}
波长 $\lambda=0.30\,\mu\mathrm m$ 的光照射
工作于饱和状态的光电管，
其光谱灵敏度为
$J=4.8\,\mathrm{mA/W}$。
求每个入射光子平均产生的光电子数（量子产额）。
\exsolution
设入射功率为 $P$，
每秒入射光子数 $P/(hc/\lambda)$。
设量子产额为 $w$，饱和光电流
$I=e wP/(hc/\lambda)$。
因此
\[
 \boxed{w=\frac{Jhc}{e\lambda}}.
\]
代 $J=4.8\times10^{-3}\,\mathrm{A/W}$、
$\lambda=3.0\times10^{-7}\,\mathrm m$，得
\[
 \boxed{w\simeq1.98\times10^{-2}\approx2.0\%}.
\]
\excheck 饱和电流时每个逸出电子均被收集，
这是用测得的电流推求量子产额的条件。

\thermalopticsproblem{278}{230}
真空光电管两极板分别由铯与铜制成，
外电路短接。用波长
$\lambda=0.22\,\mu\mathrm m$ 的光照射铯极，
求到达铜极的光电子最大速度。
\exsolution
铯与铜的逸出功在原书附录22中分别为
$A_{\mathrm{Cs}}=1.89\,\mathrm{eV}$、
$A_{\mathrm{Cu}}=4.47\,\mathrm{eV}$。
电子自铯表面出射最大动能为
$K_{\mathrm{Cs}}=hc/\lambda-A_{\mathrm{Cs}}$。
由于两极板\emph{电路短接}，
接触电势使从铯到铜的电子能量变化
为 $A_{\mathrm{Cs}}-A_{\mathrm{Cu}}$，
所以铜极前的最大动能为
\[
 K_{\mathrm{Cu}}=
 K_{\mathrm{Cs}}+A_{\mathrm{Cs}}-A_{\mathrm{Cu}}
 =\frac{hc}{\lambda}-A_{\mathrm{Cu}}.
\]
因此
\[
 \boxed{v_{\max}=\sqrt{\frac{2}{m_e}
 \left(\frac{hc}{\lambda}-A_{\mathrm{Cu}}\right)}
 \simeq6.4\times10^5\,\mathrm{m/s}}.
\]
\excheck 虽由铯产生光电子，
短路两电极间的接触势差不能忽略；
最终能量由\emph{铜}的逸出功控制。

\thermalopticsproblem{279}{230}
真空光电管锌阴极受到波长
$\lambda=262\,\mathrm{nm}$ 的辐射，
当外加减速电压 $U_{\mathrm{ext}}=1.5\,\mathrm V$
时光电流刚被阻断。
求光电管的接触电势差大小及其相对外电压的极性。
\exsolution
锌逸出功为 $A_{\mathrm{Zn}}=3.74\,\mathrm{eV}$，
由光电方程
\[
 K_{\max}=\frac{hc}{262\,\mathrm{nm}}-A_{\mathrm{Zn}}
 \simeq0.993\,\mathrm{eV}.
\]
电子实际遇到的总反向电势差的大小应为
$U_{\mathrm{stop}}=K_{\max}/e\simeq0.993\,\mathrm V$。
外加反向电压为 $1.5\,\mathrm V$，故内部接触势
必须抵消其一部分：
\[
 \boxed{|U_{\mathrm c}|=
 U_{\mathrm{ext}}-U_{\mathrm{stop}}
 \simeq0.51\,\mathrm V\approx0.5\,\mathrm V}.
\]
\excheck 接触电势的极性与\emph{外加减速电压相反}，
否则光电流会在低于 $1.5\,\mathrm V$ 的外电压下被阻断。

\thermalopticsproblem{280}{230}
利用光速 $c$、粒子质量 $m$ 和普朗克常数 $h$，
构造具有长度量纲的物理量，并说明其名称。
\exsolution
若长度尺度为 $L=h^\alpha m^\beta c^\gamma$，
依 SI 量纲
$[h]=ML^2T^{-1}$、
$[m]=M$、$[c]=LT^{-1}$，
要求 $M^0L^1T^0$，得到
$\alpha=1,\beta=-1,\gamma=-1$。
故
\[
 \boxed{\lambda_C=\frac{h}{mc}},
\]
它是质量为 $m$ 的粒子的\emph{康普顿波长}。
\excheck 此处使用普通普朗克常数 $h$；
约化康普顿波长则是 $\hbar/(mc)=\lambda_C/(2\pi)$。
% === END 33-thermal-radiation-269-280.tex ===

% === BEGIN 33-thermal-radiation-281-292.tex blob 26f6cf107fd3b0a1961384bf7e561389c3bdc89c ===
% Irodov Part V §5.7, originals 281--292 (end of Part V).
% Supplied Arihant PDF pp.240 / printed p.231; Appendix and answers pp.347-348.
% In problem 292, photon gains energy because the incident electron STOPS:
% the signed wavelength displacement is negative; book gives its magnitude.
\thermalopticsproblem{281}{231}
应用能量、动量守恒证明：一个在真空中自由运动的
电子不可能完全吸收一个真实光子。
\exsolution
在电子吸收光子前的静止参考系，
电子四动量为 $P^\mu=(m_ec,\boldsymbol0)$，
光子四动量为 $K^\mu=(E/c,\boldsymbol k)$，
其中 $E=pc>0$、$K^\mu K_\mu=0$。
若完全吸收，末态仅一个电子，其四动量为
$P_f^\mu=P^\mu+K^\mu$，应满足电子壳条件
$P_f^\mu P_{f\mu}=m_e^2c^2$。
但直接平方可得
\[
 P_f^2=(P+K)^2=P^2+2P\!\cdot\! K+K^2
 =m_e^2c^2+2m_eE>m_e^2c^2,
\]
与自由电子质量不变相矛盾。
故
\[
 \boxed{\text{自由电子不能单独完全吸收真实光子}}.
\]
\excheck 若有晶格、原子核或其他粒子共同承担反冲，
则不受此两体末态限制；本证明使用自由电子这一必要条件。

\thermalopticsproblem{282}{231}
解释康普顿散射的三个现象：
(a) 散射波长增加量与材料种类无关；
(b) 散射角越大、原子序数越低，
移位谱线的强度通常越高；
(c) 散射光中为何还保有未发生位移的谱线？
\exsolution
(a) 对与\emph{近似自由}电子的一次康普顿散射，
能量、动量守恒给
\[
 \boxed{\Delta\lambda=
 \frac{h}{m_ec}(1-\cos\theta)}.
\]
其中仅含电子质量、常数以及散射角，
故不直接依赖散射材料。
(b) 在一定能区内，高散射角意味着更大的能量转移，
更容易使结合电子表现出准自由电子行为；
低原子序数材料中电子结合较弱，
因此可观察到的移位分量比例往往较大。
(c) 被原子强束缚的电子或整个原子／晶格参与
弹性（相干）散射，较大的反冲体系质量使
波长变化很小，形成未位移的瑞利型散射分量。
\excheck (b) 涉及实验能区与电子结构，不能由
康普顿单电子公式孤立地推出普适的单调强度定律。

\thermalopticsproblem{283}{231}
一束窄单色 X 射线在材料中由自由电子散射。
观测散射角为
$\theta_1=60^\circ$ 与 $\theta_2=120^\circ$
的两束散射光波长之比为
$\eta=\lambda_2/\lambda_1=2.0$。
求原入射波长。
\exsolution
记电子康普顿波长
$\lambda_C=h/(m_ec)=2.426\,\mathrm{pm}$。
由康普顿公式，
\[
 \lambda_1=\lambda_0+\lambda_C(1-\cos60^\circ)
 =\lambda_0+\frac12\lambda_C,
\]
\[
 \lambda_2=\lambda_0+\lambda_C(1-\cos120^\circ)
 =\lambda_0+\frac32\lambda_C.
\]
由
$(\lambda_0+\tfrac32\lambda_C)/
(\lambda_0+\tfrac12\lambda_C)=2$
得
\[
 \boxed{\lambda_0=\frac{\lambda_C}{2}
 \simeq1.21\,\mathrm{pm}}.
\]
\excheck 大散射角对应较长波长，
故 $\lambda_2>\lambda_1$ 与题设 $\eta=2$ 一致。

\thermalopticsproblem{284}{231}
能量 $E_0=1.00\,\mathrm{MeV}$ 的光子
被静止自由电子康普顿散射，散射后
光子波长增加了入射波长的
$\eta=25\%$。求反冲电子动能。
\exsolution
光子的能量反比于波长，
$\lambda'=(1+\eta)\lambda_0$ 意味着
$E'=E_0/(1+\eta)$。
自由电子初始静止，能量守恒给电子动能
$T_e=E_0-E'$，从而
\[
 \boxed{T_e=E_0\frac{\eta}{1+\eta}
 =1.00\,\mathrm{MeV}\frac{0.25}{1.25}
 =0.20\,\mathrm{MeV}}.
\]
\excheck “波长增大25%”不等于
“光子能量减少25%”，因为二者反比。

\thermalopticsproblem{285}{231}
波长 $\lambda=6.0\,\mathrm{pm}$ 的光子被
一个静止自由电子以 $90^\circ$ 散射。
求：(a) 散射光子角频率；
(b) 电子反冲动能。
\exsolution
当 $\theta=90^\circ$，
康普顿式给 $\lambda'=\lambda+\lambda_C$，
$\lambda_C=2.426\,\mathrm{pm}$。
所以
\[
 \boxed{\omega'=\frac{2\pi c}{\lambda+\lambda_C}
 \simeq2.24\times10^{20}\,\mathrm{s^{-1}}}.
\]
由能量守恒，
\[
 \boxed{T_e=hc\left(\frac1\lambda-
 \frac1{\lambda+\lambda_C}\right)
 \simeq59.5\,\mathrm{keV}\approx60\,\mathrm{keV}}.
\]
\excheck (a) 求的是\emph{角频率}而非 $\nu=c/\lambda'$，
二者相差 $2\pi$。

\thermalopticsproblem{286}{231}
初能量为 $E=250\,\mathrm{keV}$ 的光子，
被静止自由电子以角 $\theta=120^\circ$ 散射。
求散射光子的能量。
\exsolution
康普顿公式
$\lambda'-\lambda=\lambda_C(1-\cos\theta)$，
以 $E=hc/\lambda$ 改写为能量形式
\[
 E'=\frac E{1+
 \dfrac E{m_ec^2}(1-\cos\theta)}.
\]
由于
$1-\cos120^\circ=3/2$、
$m_ec^2=511\,\mathrm{keV}$，
\[
 \boxed{E'=
 \frac{250}{1+(250/511)\times1.5}\,\mathrm{keV}
 \simeq144\,\mathrm{keV}}.
\]
\excheck $E'<E$，且散射角较大时
转移给电子的能量较多。

\thermalopticsproblem{287}{231}
初始光子动量为
$p=1.02\,\mathrm{MeV}/c$，
散射后动量变为
$p'=0.255\,\mathrm{MeV}/c$，
散射体为静止自由电子。
求光子的散射角。
\exsolution
因 $\lambda=h/p$、
$\lambda'=h/p'$，
康普顿位移方程为
\[
 \frac hp'-\frac hp=
 \frac h{m_ec}(1-\cos\theta).
\]
约去 $h$ 得
\[
 \boxed{1-\cos\theta=
 m_ec\left(\frac1{p'}-\frac1p\right)}.
\]
以 $m_ec^2=0.511\,\mathrm{MeV}$ 代入，
\[
 1-\cos\theta=
 0.511\left(\frac1{0.255}-\frac1{1.02}\right)
 \simeq1.503,
\]
数值取有效位数后 $\cos\theta\simeq-0.50$，
故
\[
 \boxed{\theta\simeq120^\circ}.
\]
\excheck 使用 $\mathrm{MeV}/c$ 表达动量时，
$m_ec$ 也须换成相同的动量单位。

\thermalopticsproblem{288}{231}
光子被静止自由电子以
$\theta=120^\circ$ 散射后，
电子动能为 $T=0.45\,\mathrm{MeV}$。
求散射前光子的能量 $E$。
\exsolution
由能量守恒，散射光子能量 $E'=E-T$。
康普顿能量关系
\[
 \frac1{E'}-\frac1E=
 \frac{1-\cos\theta}{m_ec^2}.
\]
于是
\[
 \frac{T}{E(E-T)}
 =\frac{1-\cos\theta}{m_ec^2}.
\]
令 $A=m_ec^2/(1-\cos\theta)$，得
$E^2-TE-AT=0$，取正根：
\[
 \boxed{E=\frac T2
 \left[1+\sqrt{1+
 \frac{4m_ec^2}{T(1-\cos\theta)}}\right]}.
\]
代入 $T=0.45\,\mathrm{MeV}$、
$1-\cos120^\circ=1.5$，
\[
 \boxed{E\simeq0.68\,\mathrm{MeV}}.
\]
\excheck 解的另一根为负，不满足
$E>T>0$ 的物理条件。

\thermalopticsproblem{289}{231}
在自由电子康普顿散射中，
反冲电子的\emph{最大}动能为
$T_{\max}=0.19\,\mathrm{MeV}$。
求入射 X 射线波长。
\exsolution
电子得到最大能量当光子背散射 $\theta=\pi$。
若入射能量为 $E$，散射后能量为
$E'=E/(1+2E/(m_ec^2))$，
所以
\[
 T_{\max}=E-E'=
 \frac{2E^2}{m_ec^2+2E}.
\]
解二次方程得正根
\[
 \boxed{E=
 \frac{T_{\max}+
 \sqrt{T_{\max}^2+2m_ec^2T_{\max}}}{2}}.
\]
随后 $\lambda=hc/E$，数值为
\[
 \boxed{\lambda\simeq3.7\,\mathrm{pm}}.
\]
\excheck 最大反冲对应 $\theta=180^\circ$，
不是 $90^\circ$。

\thermalopticsproblem{290}{231}
能量 $E=0.15\,\mathrm{MeV}$ 的光子
被静止自由电子散射后，波长增加
$\Delta\lambda=3.0\,\mathrm{pm}$。
求反冲\emph{电子}与入射方向夹角 $\phi$。
\exsolution
光子散射角 $\theta$ 先由
$\Delta\lambda=\lambda_C(1-\cos\theta)$ 求出，
$\cos\theta=1-3.0/2.426\simeq-0.237$。
令初、末光子能量为 $E,E'$，
以能量为动量单位 $pc=E$，
$E'=hc/(\lambda+\Delta\lambda)$，
$\lambda=hc/E$，得
$E'\simeq0.110\,\mathrm{MeV}$。
动量守恒的平行、垂直入射方向分量分别为
\[
 p_{e,x}=\frac{E-E'\cos\theta}{c},
 \qquad p_{e,y}=-\frac{E'\sin\theta}{c}.
\]
所以
\[
 \boxed{\tan\phi=
 \frac{E'\sin\theta}{E-E'\cos\theta}},
 \qquad
 \boxed{\phi\simeq31^\circ}.
\]
\excheck 题目求的是\emph{电子}方向，
不等于光子散射角 $\theta\simeq104^\circ$。

\thermalopticsproblem{291}{231}
能量为电子静能的 $\eta=2.0$ 倍的光子，
与静止自由电子正碰并背向散射。
反冲电子随后以速度垂直于均匀磁场
$B=0.12\,\mathrm T$ 运动。
求其轨道曲率半径。
\exsolution
记 $m=m_e$，
$E=\eta mc^2$。
光子完全背散射时
\[
 E'=\frac{E}{1+2\eta},\qquad
 T=E-E'=\frac{2\eta^2}{1+2\eta}mc^2.
\]
电子满足相对论动量--能量关系
$p_ec=\sqrt{T(T+2mc^2)}$。
代入并化简得
\[
 p_e=\frac{2\eta(1+\eta)}{1+2\eta}mc.
\]
垂直磁场中洛伦兹力提供向心加速度，
相对论曲率半径为 $r=p_e/(eB)$，所以
\[
 \boxed{r=\frac{2\eta(1+\eta)m_ec}{(1+2\eta)eB}
 \simeq3.4\,\mathrm{cm}}.
\]
\excheck 应使用\emph{相对论动量}，
不可在 MeV 电子上直接代 $p=mv_{\rm nonrel}$。

\thermalopticsproblem{292}{231}
一个运动的相对论电子与光子碰撞。
光子方向偏转 $\theta=60^\circ$，
而电子在碰撞后\emph{恰好停止}。
求散射光子的康普顿波长变化。
\exsolution
与通常“静止电子受撞击”不同，
此题入射电子在运动而末态电子静止。
令入射、出射光子四动量为 $K,K'$，
末态电子四动量为 $P_f=(m_ec,\boldsymbol0)$。
动量守恒写作
$P_i+K=P_f+K'$；
利用 $P_i^2=P_f^2=m_e^2c^2$、
$K^2=K'^2=0$，
平方后得到
\[
 m_ec^2(E'-E)=EE'(1-\cos\theta).
\]
所以
$\tfrac1E-\tfrac1{E'}=(1-\cos\theta)/(m_ec^2)$。
换成波长，得到带符号的光子波长变化
\[
 \boxed{\lambda'-\lambda=
 -\frac h{m_ec}(1-\cos\theta)}.
\]
在 $\theta=60^\circ$，
\[
 \boxed{\Delta\lambda\simeq-1.21\,\mathrm{pm}},
 \qquad
 \boxed{|\Delta\lambda|\simeq1.2\,\mathrm{pm}}.
\]
\excheck 此题是\emph{反向康普顿能量传递}，
电子停止使光子能量增加、波长缩短；
书末只给波长变化的模，负号须独立标注。
% === END 33-thermal-radiation-281-292.tex ===
